% Source: https://artofproblemsolving.com/community/c3h1365819_1984_aime_flawed_solution
% Stable identity: aops-topic-1365819-q04
Let $S$ have $n$ numbers. Their sum is $56n$, and after removing one $68$ the remaining $n-1$ numbers have sum $55(n-1)$. So
\[
56n-68=55(n-1)\quad\Longrightarrow\quad n=13 .
\]
The sum of all $13$ numbers is $56\cdot13=728$, and the $12$ numbers other than the removed $68$ have sum $728-68=660$.

To make one of these $12$ numbers as large as possible, make the other $11$ as small as possible, namely all equal to $1$. Then the largest number is
\[
660-11\cdot1=\boxed{649},
\]
and the list ($68$, $649$ and eleven $1$s) works: its average is $\dfrac{728}{13}=56$, and without $68$ it is $\dfrac{660}{12}=55$.

The value $659$ would leave only one other number besides $68$ and the largest one, but there are $11$ other numbers, each at least $1$. So $649$ is right and the wiki solution is correct.
